This research note contributes to the literature on family and pregnant women's health and well-being household structure and women’s well being in India by revisiting the relationship between household structure, autonomy, and maternal healthcare among pregnant women during a period in which earlier trends toward nuclearization appear to have reversed. we show that despite the growing prevalence of patrilocal joint residence, pregnant women's autonomy indicators have improved both overall, and within each household structure. 

These contrasting dynamics suggest that population level changes in household structure can have ambiguous impacts on maternal health - although newly married women in patrilocal households may benefit from pooled resources and familial support, they are also disadvantaged by their low intrahousehold status. Understanding whether shifts in household structure affects maternal healthcare and women’s autonomy is critical for interpreting trends in women’s well being in contemporary India. 

Across the survey rounds, women’s autonomy and maternal healthcare utilization improved steadily, alongside changes in the distribution of pregnant women across household structures. To distinguish between improvements in these outcomes which occurred within household structures and improvements attributable to changing composition across household structures, we use Kitagawa decomposition.  The decomposition results indicate that most improvements in autonomy and maternal healthcare reflect changes within household types, with only a small share attributable to compositional change. 



More broadly, our findings suggest that household structure shapes different dimensions of women’s well being through distinct processes. The higher maternal healthcare use observed in patrilocal joint households is closely tied to differences in economic resources, whereas  in women’s decision making power persist beyond wealth. These patterns indicate that improvements in maternal conditions do not necessarily weaken intra household hierarchies. Taken together the results show the importance of analyzing household structure not as a shift from ‘traditional’ to ‘modern’ but as a set of family arrangements that differently shape women’s autonomy and access to care. 